\documentclass[../DoAn.tex]{subfiles}
\begin{document}

\section{Giới thiệu chương}

Chương này trình bày những giải pháp và đóng góp nổi bật mà em đã xây dựng trong
quá trình thực hiện đề tài hệ thống giám sát và bảo vệ xe. Thay vì mô tả lại
toàn bộ chức năng đã được giới thiệu ở các chương trước, nội dung chương tập
trung vào các vấn đề có tính kỹ thuật, những khó khăn phát sinh trong quá trình
thiết kế, các phương án đã được cân nhắc, giải pháp được lựa chọn và kết quả đạt
được.
  
Bốn đóng góp chính được trình bày gồm: giải thuật phát hiện xe di chuyển bất
thường; cơ chế cảnh báo đa kênh kết hợp Socket.IO và Firebase Cloud Messaging;
phương pháp tái tạo hành trình và tính quãng đường từ dữ liệu GPS; giải pháp
quản lý vòng đời và quyền điều khiển thiết bị IoT. Đây là các nội dung thể hiện
rõ nhất quá trình phân tích, giải quyết vấn đề và hoàn thiện hệ thống của tác
giả.


\section{Giải thuật phát hiện xe di chuyển bất thường}
\label{sec:phat-hien-bat-thuong}

\subsection{Bài toán và khó khăn}

Mục tiêu quan trọng nhất của hệ thống là phát hiện trường hợp xe bị di chuyển
khi người dùng đã bật chế độ bảo vệ. Nếu chỉ kiểm tra xem xe có tốc độ lớn hơn
không, hệ thống có thể bỏ sót trường hợp dữ liệu tốc độ chưa được cập nhật nhưng
tọa độ GPS đã thay đổi. Ngược lại, nếu chỉ dựa trên khoảng cách giữa hai tọa độ,
sai số tự nhiên của GPS có thể làm hệ thống hiểu nhầm rằng xe đang di chuyển dù
xe vẫn đứng yên.

Dữ liệu GPS còn có thể xuất hiện các điểm trôi, điểm nhảy hoặc bản tin không ổn
định do chất lượng tín hiệu vệ tinh. Một bản tin bất thường đơn lẻ vì vậy chưa
đủ độ tin cậy để kết luận xe đang gặp nguy hiểm. Việc phát cảnh báo ngay từ bản
tin đầu tiên sẽ làm tăng tỷ lệ cảnh báo giả, gây khó chịu và khiến người dùng
dần bỏ qua các cảnh báo thật.

Ngoài ra, cùng một dữ liệu tốc độ và vị trí phải được xử lý khác nhau tùy theo
trạng thái của xe. Khi người dùng tắt bảo vệ, xe di chuyển là hành vi bình
thường. Khi người dùng bật bảo vệ, vị trí tại thời điểm đó trở thành vị trí tham
chiếu và mọi dịch chuyển đáng kể sau đó cần được kiểm tra. Vì vậy, bài toán cần
kết hợp trạng thái bảo vệ, tốc độ, khoảng cách và lịch sử các bản tin liên tiếp.

\subsection{Giải pháp đề xuất}

Khi chế độ bảo vệ được bật, hệ thống ghi nhận vị trí gần nhất của xe làm vị trí
tham chiếu. Với mỗi bản tin MQTT mới, backend chuẩn hóa tốc độ và tọa độ trước
khi đánh giá rủi ro. Khoảng cách giữa vị trí hiện tại và vị trí bảo vệ được tính
bằng công thức Haversine.

Gọi $\varphi_1$, $\varphi_2$ là vĩ độ; $\lambda_1$, $\lambda_2$ là kinh độ của
hai điểm tính theo radian. Khoảng cách được xác định như sau:

\begin{equation}
 a = \sin^2\left(\frac{\Delta\varphi}{2}\right)
 + \cos(\varphi_1)\cos(\varphi_2)
 \sin^2\left(\frac{\Delta\lambda}{2}\right)
\end{equation}

\begin{equation}
 d = 2R\arcsin(\sqrt{a})
\end{equation}

Trong đó, $R = 6.371.000$ m là bán kính trung bình của Trái Đất. Hệ thống sử
dụng các điều kiện đánh giá sau:
Bất thường tốc độ được xác định khi tốc độ lớn hơn 10 km/h. Bất thường khoảng cách được xác định khi khoảng cách di chuyển lớn hơn 50 m. Trường hợp di chuyển nghiêm trọng xảy ra khi khoảng cách di chuyển lớn hơn 100 m. Một bản ghi được coi là dữ liệu nghi ngờ khi xuất hiện bất thường tốc độ hoặc bất thường khoảng cách.

Để giảm cảnh báo giả, biến \texttt{suspiciousCount} lưu số bản tin bất thường
liên tiếp. Bản tin bất thường đầu tiên chỉ làm tăng bộ đếm. Cảnh báo được tạo khi
hệ thống nhận đủ hai bản tin bất thường liên tiếp. Nếu dữ liệu trở lại bình
thường hoặc người dùng tắt chế độ bảo vệ, bộ đếm được đặt lại. Biến
\texttt{lastSpeedAlert} được sử dụng để ngăn gửi nhiều cảnh báo cho cùng một
đợt bất thường.

Xe được đánh dấu ở trạng thái nguy hiểm khi khoảng cách vượt 100 m hoặc đồng
thời xuất hiện cả bất thường về tốc độ và khoảng cách. Quy trình xử lý tổng quát
được mô tả như sau:

\begin{itemize}
\item Nhận dữ liệu từ thiết bị thông qua giao thức MQTT.
\item Kiểm tra thông tin thiết bị và trạng thái kích hoạt.
\item Lưu dữ liệu cảm biến vào cơ sở dữ liệu.
\item Kiểm tra chế độ bảo vệ của phương tiện.
\item Chuẩn hóa tốc độ và tọa độ.
\item Tính khoảng cách di chuyển bằng công thức Haversine.
\item Kiểm tra các dấu hiệu bất thường.
\item Cập nhật bộ đếm bất thường liên tiếp.
\item Tạo cảnh báo nếu thỏa mãn điều kiện cảnh báo.
\end{itemize}

Giải pháp này có thể tổng quát hóa cho các hệ thống giám sát đối tượng di động
khác bằng cách thay đổi ngưỡng tốc độ, khoảng cách, số mẫu liên tiếp và quy tắc
phân loại mức độ nghiêm trọng.

\subsection{Kết quả đạt được và đánh giá}

Giải thuật đã được tích hợp vào luồng tiếp nhận dữ liệu MQTT của backend. Dữ
liệu cảm biến được lưu trước khi đánh giá rủi ro, nhờ đó có thể đối chiếu lại
lịch sử khi cần. Khi đủ điều kiện, hệ thống tạo bản ghi cảnh báo, cập nhật trạng
thái thiết bị và chuyển thông tin đến ứng dụng.

\begin{table}[ht]
\centering
\label{tab:ket-qua-phat-hien}
\begin{tabular}{|p{0.48\textwidth}|p{0.40\textwidth}|}
\hline
\textbf{Tình huống} & \textbf{Kết quả} \\ \hline
Tốc độ không vượt 10 km/h và khoảng cách không vượt 50 m & Không tạo cảnh báo \\ \hline
Chỉ có một bản tin vượt ngưỡng & Tăng bộ đếm, chưa tạo cảnh báo \\ \hline
Hai bản tin liên tiếp vượt ngưỡng & Tạo cảnh báo mới \\ \hline
Khoảng cách vượt 100 m & Tạo cảnh báo khẩn cấp và đánh dấu nguy hiểm \\ \hline
Đồng thời vượt tốc độ và khoảng cách & Tạo cảnh báo khẩn cấp \\ \hline
Dữ liệu trở lại bình thường & Đặt lại bộ đếm bất thường \\ \hline
\end{tabular}
\caption{Kết quả mong đợi của giải thuật phát hiện bất thường}
\end{table}

Ưu điểm của giải pháp là không phụ thuộc vào một nguồn thông tin duy nhất và có
cơ chế chống cảnh báo lặp. Tuy nhiên, các ngưỡng hiện được lựa chọn dựa trên môi
trường thử nghiệm và chưa được hiệu chỉnh bằng tập dữ liệu thực tế đủ lớn.


\section{Cơ chế cảnh báo đa kênh}
\label{sec:canh-bao-da-kenh}

\subsection{Bài toán và khó khăn}

Ứng dụng cần thông báo cho người dùng trong nhiều trạng thái khác nhau. Khi ứng
dụng đang mở và kết nối mạng ổn định, dữ liệu cần được cập nhật gần thời gian
thực. Khi ứng dụng chạy nền hoặc đã bị hệ điều hành tạm dừng, kết nối Socket.IO
có thể không còn hoạt động. Nếu chỉ sử dụng Socket.IO, người dùng có nguy cơ
không nhận được cảnh báo quan trọng. Nếu chỉ sử dụng push notification, giao
diện khi đang mở có thể cập nhật chậm và không tạo được cảm giác realtime.

Một khó khăn khác là dữ liệu chỉ được phép gửi đến đúng chủ sở hữu thiết bị.
Backend không thể phát toàn bộ cảnh báo cho tất cả kết nối đang hoạt động. Đồng
thời, khi cảnh báo được nhận, trạng thái của xe trên trang chủ, danh sách thiết
bị, màn hình chi tiết và danh sách cảnh báo phải được đồng bộ mà không yêu cầu
người dùng tải lại trang.

\subsection{Giải pháp đề xuất}

Hệ thống sử dụng hai kênh cảnh báo bổ trợ cho nhau. Socket.IO phục vụ cập nhật
realtime khi ứng dụng đang hoạt động. Firebase Cloud Messaging đóng vai trò
kênh thông báo khi ứng dụng chạy nền, bị tạm dừng hoặc kết nối Socket.IO không
sẵn sàng.

Quá trình gửi cảnh báo được thực hiện bởi `SensorDataService` gồm các bước sau:

\begin{itemize}
\item Lưu thông tin cảnh báo vào cơ sở dữ liệu PostgreSQL.
\item Gửi sự kiện cảnh báo thời gian thực thông qua `SensorDataGateway` sử dụng Socket.IO.
\item Gửi thông báo đẩy đến người dùng thông qua `FirebaseMessagingService` và Firebase Cloud Messaging.
\end{itemize}

Mỗi client Socket.IO gửi JWT sau khi kết nối. Backend xác thực token, lấy
\texttt{userId} và gắn định danh người dùng với socket tương ứng. Khi phát sinh
cảnh báo, gateway chỉ gửi sự kiện đến các socket có cùng \texttt{userId}. Đối
với FCM, token của từng thiết bị di động được lưu trong bảng
\texttt{FcmToken}. Backend lấy toàn bộ token hợp lệ của người dùng và gửi cảnh
báo multicast.

Ở phía Flutter, sự kiện realtime không chỉ được chuyển cho bộ điều khiển cảnh
báo mà còn được chuyển cho bộ điều khiển thiết bị. Trường
\texttt{vehicleStatus} trong payload được dùng để cập nhật ngay trạng thái
\texttt{Stolen}. Vì vậy, trang chủ, danh sách thiết bị và màn hình chi tiết cùng
thay đổi mà không cần gọi lại API. Khi người dùng mở ứng dụng từ FCM, hệ thống
tải lại thiết bị và cảnh báo để đồng bộ với database.

Thông báo vẫn được nhận dù người dùng đang lọc danh sách theo ngày hoặc thiết
bị. Bộ lọc chỉ giới hạn dữ liệu hiển thị trong danh sách, không ngăn local
notification hoặc push notification khẩn cấp.

\subsection{Kết quả đạt được và đánh giá}

Giải pháp giúp ứng dụng hoạt động phù hợp ở cả foreground và background. Khi
ứng dụng đang mở, Socket.IO cập nhật giao diện trực tiếp. Khi ứng dụng chạy nền,
FCM tạo thông báo hệ thống và điều hướng người dùng đến màn hình cảnh báo khi
được nhấn.

Backend đồng thời loại bỏ các FCM token đã hết hiệu lực dựa trên phản hồi của
Firebase. Điều này hạn chế việc tiếp tục gửi đến token không còn tồn tại. Cảnh
báo được lưu vào database trước khi gửi qua hai kênh, nhờ đó người dùng vẫn có
thể xem lại lịch sử nếu việc gửi realtime hoặc push gặp sự cố.

Hạn chế của giải pháp là chưa xây dựng cơ chế hàng đợi và thử gửi lại khi dịch
vụ ngoài tạm thời không khả dụng. Hệ thống hiện cũng chưa lưu trạng thái đã đọc
của từng cảnh báo. Các nội dung này có thể được phát triển bằng message queue,
retry có giới hạn và trường \texttt{readAt} trong bảng cảnh báo.

\section{Tái tạo hành trình và tính quãng đường từ dữ liệu GPS}
\label{sec:tai-tao-hanh-trinh}

\subsection{Bài toán và khó khăn}

Dữ liệu GPS được thiết bị gửi theo thời gian, nhưng không phải mọi điểm đều có
thể sử dụng trực tiếp để vẽ hành trình. Tọa độ có thể sử dụng các tên trường
khác nhau như \texttt{lat}, \texttt{latitude}, \texttt{lng} hoặc
\texttt{longitude}. Một số bản ghi có thể thiếu tọa độ, vượt phạm vi hợp lệ,
trùng nhau hoặc xuất hiện bước nhảy lớn do mất tín hiệu.

Nếu cộng khoảng cách giữa mọi cặp điểm mà không lọc, sai số vài mét khi xe đứng
yên sẽ tích lũy thành một quãng đường đáng kể. Bước nhảy GPS cũng có thể làm
tổng quãng đường tăng hàng kilomet dù xe chỉ di chuyển trong phạm vi nhỏ. Vì
vậy, chức năng này cần vừa tái tạo đúng thứ tự hành trình vừa hạn chế nhiễu.

\subsection{Giải pháp đề xuất}

Backend nhận khoảng thời gian cần truy vấn, lấy các bản ghi cảm biến theo thứ tự
tăng dần của \texttt{createdAt} và chuẩn hóa tọa độ về cấu trúc thống nhất gồm
\texttt{lat} và \texttt{lng}. Các điểm nằm ngoài phạm vi vĩ độ $[-90,90]$ hoặc
kinh độ $[-180,180]$ bị loại bỏ.

Khoảng cách giữa điểm mới và điểm hợp lệ gần nhất được tính bằng Haversine. Một
điểm bị bỏ qua nếu khoảng cách nhỏ hơn 3 m vì nhiều khả năng đây là rung GPS khi
xe đứng yên. Hệ thống cũng tính vận tốc suy ra từ khoảng cách và chênh lệch thời
gian. Nếu vận tốc suy ra vượt 200 km/h, điểm mới được xem là bước nhảy không hợp
lý và không được cộng vào tổng quãng đường.

Quy trình tính toán quãng đường di chuyển của thiết bị gồm các bước sau:

\begin{itemize}
\item Lấy danh sách các điểm dữ liệu theo thứ tự thời gian tăng dần.
\item Chuẩn hóa tọa độ của từng điểm dữ liệu.
\item Loại bỏ các tọa độ không hợp lệ.
\item Loại bỏ các điểm có mức dịch chuyển nhỏ hơn 3 m để giảm nhiễu GPS.
\item Loại bỏ các điểm có vận tốc suy ra lớn hơn 200 km/h do có khả năng là dữ liệu sai lệch.
\item Tính và cộng dồn khoảng cách giữa các điểm liên tiếp bằng công thức Haversine.
\item Trả về danh sách các điểm hợp lệ và tổng quãng đường di chuyển.
\end{itemize}

App sẽ chuyển các điểm hợp lệ thành danh sách \texttt{LatLng} và sử dụng
\texttt{PolylineLayer} để vẽ hành trình trên OpenStreetMap. Điểm đầu, điểm cuối,
số điểm GPS và tổng quãng đường được hiển thị trực tiếp. Người dùng có thể chọn
phạm vi 24 giờ, 7 ngày hoặc toàn bộ lịch sử.

\subsection{Kết quả đạt được và đánh giá}

Chức năng đã cho phép người dùng quan sát tuyến đường của xe thay vì chỉ xem vị
trí cuối cùng. Tổng quãng đường được tính ở backend nên kết quả không phụ thuộc
vào khả năng xử lý của điện thoại. Việc lọc theo thời gian cũng làm giảm lượng
dữ liệu truyền đến ứng dụng.

Các quy tắc 3 m và 200 km/h giúp loại bỏ các trường hợp nhiễu rõ ràng, nhưng
chưa thay thế được các thuật toán làm mượt GPS chuyên dụng. Khi có dữ liệu thực
tế, có thể đánh giá thêm bộ lọc Kalman, map matching hoặc sử dụng độ chính xác
do module GPS cung cấp. Đây là hướng tổng quát hóa giải pháp cho các hệ thống
theo dõi phương tiện có yêu cầu độ chính xác cao hơn.

\section{Quản lý vòng đời và quyền điều khiển thiết bị IoT}
\label{sec:quan-ly-thiet-bi-iot}

\subsection{Bài toán và khó khăn}

Thiết bị IoT được tạo trước khi thuộc về một tài khoản cụ thể. Nếu chỉ sử dụng
\texttt{deviceId}, một người biết mã thiết bị có thể kích hoạt thiết bị không
thuộc sở hữu của mình. Hệ thống cũng phải ngăn người dùng xem dữ liệu, thay đổi
trạng thái bảo vệ hoặc điều khiển còi của thiết bị thuộc người dùng khác.

Ở chiều MQTT, backend nhận dữ liệu từ broker thay vì nhận trực tiếp từ một kết
nối HTTP đã có JWT. Do đó cần một cơ chế riêng để xác thực nguồn gửi dữ liệu mà
không làm phát sinh truy vấn database cho mọi bản tin cảm biến.

\subsection{Giải pháp đề xuất}

Mỗi thiết bị được định danh bằng \texttt{deviceId} và một khóa bí mật
\texttt{deviceKey}. Khi kích hoạt, người dùng phải đăng nhập và cung cấp đúng cả
hai giá trị. Backend kiểm tra thiết bị chưa thuộc người dùng khác, sau đó cập
nhật \texttt{userId}, \texttt{isActivated} và \texttt{activatedAt}.

Các API dành cho người dùng đều kiểm tra JWT. Sau khi tìm thiết bị, service so
sánh \texttt{device.userId} với \texttt{req.user.id}. Yêu cầu không đúng chủ sở
hữu bị từ chối trước khi truy vấn dữ liệu hoặc gửi lệnh MQTT.

Thiết bị gửi \texttt{deviceKey} trong payload MQTT. Kết quả
xác thực thành công được lưu trong cache một giờ. Các bản tin tiếp theo của cùng
thiết bị không cần truy vấn database lại trong thời gian cache còn hiệu lực.

Luồng điều khiển còi cũng kiểm tra quyền sở hữu trước khi publish lệnh đến topic
\texttt{buzzer/deviceId}. Trạng thái còi trong database chỉ được cập nhật sau
khi yêu cầu đã vượt qua bước kiểm tra quyền.

\subsection{Kết quả đạt được và đánh giá}

Giải pháp bảo đảm một thiết bị không thể được kích hoạt chỉ bằng mã công khai,
không thể đồng thời thuộc nhiều người dùng và không thể bị điều khiển bởi tài
khoản không có quyền. Cache xác thực làm giảm số truy vấn database khi thiết bị
gửi dữ liệu với tần suất cao.

Tuy vậy, \texttt{deviceKey} hiện vẫn là khóa tĩnh. Firmware thử nghiệm còn chứa
thông tin cấu hình trực tiếp và sử dụng kết nối chưa kiểm tra đầy đủ chứng thư.
Khi triển khai thực tế, khóa cần được lưu trong vùng bảo mật, hỗ trợ xoay vòng,
không ghi vào mã nguồn và kết nối MQTT cần xác minh chứng thư máy chủ.

\section{Tổng kết chương}

Chương đã trình bày bốn đóng góp tập trung vào các vấn đề cốt lõi của hệ thống
giám sát và bảo vệ xe. Giải thuật phát hiện bất thường kết hợp trạng thái bảo vệ,
tốc độ, khoảng cách và chuỗi bản tin liên tiếp để giảm cảnh báo giả. Cơ chế cảnh
báo đa kênh giúp hệ thống hoạt động ở cả foreground và background, đồng thời
đồng bộ trạng thái giao diện theo thời gian thực. Phương pháp xử lý GPS cung cấp
hành trình và quãng đường có lọc nhiễu cơ bản. Giải pháp kích hoạt và phân quyền
thiết bị tạo ra ranh giới rõ ràng giữa người dùng, thiết bị và kênh MQTT.


\end{document}